\section{Results}
\label{sec:results}

All results below come from single-seed runs at one model scale. Section~\ref{sec:instrument}
describes three defects in the measurement apparatus found and corrected before these were
collected; results predating those corrections have been discarded.

\subsection{Part 1: does stop-gradient still matter under SIGReg?}

We swept the SIGReg strength $\lambda$ across the range \citet{lejepa2025} itself uses,
running both SIGReg conditions at each value for 2{,}000 steps. Mean pairwise cosine
similarity near zero indicates a healthy, spread-out representation; near one indicates
collapse.

\begin{table}[h]
\centering
\small
\begin{tabular}{r r r r r}
\toprule
& \multicolumn{2}{c}{\textbf{without stop-gradient}} & \multicolumn{2}{c}{\textbf{with stop-gradient}} \\
\cmidrule(lr){2-3} \cmidrule(lr){4-5}
$\lambda$ & cosine & total var. & cosine & total var. \\
\midrule
0.01 & \textbf{0.990} & 0.56  & 0.034 & 153.4 \\
0.02 & \textbf{0.951} & 1.83  & 0.021 & 162.5 \\
0.05 & \textbf{0.765} & 19.3  & 0.009 & 167.6 \\
0.10 & 0.381          & 89.5  & 0.008 & 171.0 \\
\bottomrule
\end{tabular}
\caption{\ours{} Final collapse metrics after 2{,}000 steps. With stop-gradient, no
configuration collapses at any strength. Without it, every configuration inside
\citet{lejepa2025}'s swept range collapses, though increasing $\lambda$ helps monotonically.}
\end{table}

\ours{} Within the hyperparameter range used by \citet{lejepa2025}, SIGReg alone did not
substitute for stop-gradient in our setup. The dose-response is clean and monotone --- cosine
falls from 0.990 to 0.381 as $\lambda$ rises --- so a stronger setting may eventually suffice.
$\lambda = 0.1$ is already borderline by our collapse criterion. That lies above the range the
reference sweeps and we have not tested it.

\interp{} The original framing of Part 1 asked whether a gap would \emph{widen with training
duration}. It did not need 32{,}000 steps to appear: the separation is unambiguous by 2{,}000.
We therefore report Part 1 on collapse metrics rather than on the duration-scaling question as
originally posed, since the effect was not marginal enough to require it.

\subsection{Part 2: do cheap diagnostics degrade before the probe?}

This question requires a run that both learns and subsequently degrades. Only
\texttt{ema\_stopgrad} did; the SIGReg conditions never exceeded their own random
initialisation, so their flat probe curves had nothing to lag behind.

\begin{table}[h]
\centering
\small
\begin{tabular}{l r r}
\toprule
\textbf{signal} & \textbf{turns at step} & \textbf{magnitude of move} \\
\midrule
unscaled probe accuracy (peak)   & 11{,}000 & 0.021 decline \\
mean pairwise cosine (trough)    & 13{,}500 & 0.236 rise \\
\bottomrule
\end{tabular}
\caption{\ours{} On the 32{,}000-step baseline the probe turns 2{,}500 steps \emph{before} the
cheap diagnostic --- the opposite of the hypothesis. The cheap diagnostic's eventual signal is
an order of magnitude larger, but arrives later.}
\end{table}

\ours{} The hypothesis is not supported by this run. The cheap diagnostic is a clearer
indicator of \emph{how severe} degradation has become, but not an earlier one.

\interp{} Confidence is low. This is one condition, one seed, and a 2{,}500-step lead is five
checkpoints at our logging resolution; the probe's decline is four times its measured
run-to-run noise floor of 0.005. A fair test of Part 2 needs several runs that genuinely learn
and then degrade.

\subsection{An incidental observation: the healthy baseline re-collapses}

\ours{} \texttt{ema\_stopgrad} reaches its best representation at step 11{,}000 --- roughly a
third of the budget --- and degrades slowly thereafter, with mean pairwise cosine rising from
0.294 to 0.530 and total variance falling sevenfold across the run. The classic
stop-gradient-plus-EMA recipe does not hold a representation indefinitely at this scale; it
partially re-collapses given long enough.

\subsection{What these results do not show}

\interp{} In our setup no SIGReg configuration produced a useful representation: every run
finished at or below its own random initialisation (approximately 0.369 probe accuracy), while
the EMA baseline reached 0.603. This must not be read as evidence against \citet{lejepa2025}.
We pair SIGReg with I-JEPA-style masked latent prediction; LeJEPA pairs it with multi-view
invariance across eight augmented views, with no predictor and no teacher-student network.
These are different systems sharing one component. The result is a statement about that
combination, not about SIGReg or about LeJEPA's claim.
